\documentclass{article}
\usepackage[UTF8]{ctex}
\usepackage{amsmath,amssymb,amsfonts}
\usepackage{enumitem}
\usepackage{geometry}
\geometry{a4paper, margin=1in}
% === 错题本排版优化（不参与题目标识） ===
\setlength{\parindent}{0pt}
\setlength{\parskip}{0.35em plus 0.10em minus 0.08em}
\setlength{\emergencystretch}{3em}
\setlist{itemsep=0.25em, topsep=0.40em, parsep=0pt}
\allowdisplaybreaks[2]
% === 错题本排版优化结束 ===
\begin{document}
\section*{东华大学2025年数学分析真题}
\section*{48．东华大学2025年数学分析试卷}

\subsection{题目1}
一．计算题．每题8分，共40分．

\begin{itemize}
  \item[（1）] 求极限 \(\lim _{n \rightarrow \infty} n \sin (2 \pi n \mathrm{le})\) 。
  \item[（2）] 求极限 \(\lim _{x \rightarrow 0}\left(\frac{\sin x}{x}\right)^{\frac{1}{1-\cos x}}\) ．
  \item[（3）] 求积分 \(\int_{0}^{\frac{\pi}{4}} \ln \sin (2 x) \mathrm{d} x\) 。
  \item[（4）] 求曲线积分 \(\oint_{L} \frac{y \mathrm{~d} x-x \mathrm{~d} y}{3 x^{2}-2 x y+3 y^{2}}\) ，其中 \(L\) 为 \(|x|+|y|=1\) ，沿逆时针一周．
  \item[（5）] 设 \(\Omega=\left\{(x, y, z) \mid x^{2}+y^{2}+z^{2} \leq 1, x \geq 0, y \geq 0, z \geq 0\right\}\) ，求 \(\iiint_{\Omega}\left(\frac{x^{2}}{a^{2}}+\frac{y^{2}}{b^{2}}+\frac{z^{2}}{c^{2}}\right) \mathrm{d} V\) ．
\end{itemize}

\subsection{题目2}
二．（15 分）设 \(f(x)\) 在区间 \([a, b]\) 上为增函数，且 \(f(a) \geq a, f(b) \leq b\) ，证明：存在 \(\xi \in[a, b]\) ，使得\(f(\xi)=\xi\) ．

\subsection{题目3}
三．（15 分）设函数 \(f(x)\) 在区间［0，2024］上连续（应该是连续可微），且 \(f(0)=0\) ．证明：

\[
\int_{0}^{2024}\left|f(x) f^{\prime}(x)\right| \mathrm{d} x \leq 1012 \int_{0}^{2024}\left(f^{\prime}(x)\right)^{2} \mathrm{~d} x
\]

\subsection{题目4}
四．（15 分）讨论函数

\[
f(x, y)= \begin{cases}x-y+\frac{x y^{3}}{x^{2}+y^{4}}, & x^{2}+y^{2} \neq 0 \\ 0, & x^{2}+y^{2}=0\end{cases}
\]

在原点 \((0,0)\) 是否连续，偏导数是否存在，是否可微？

\subsection{题目5}
五．解答如下问题．每题15分，共45分．

\begin{itemize}
  \item[（1）] 判断并证明级数 \(\sum_{n=1}^{\infty} \sin \left(\sqrt{n^{2}+1} \pi\right)\) 的收敛性。
\end{itemize}

\subsection{题目4}
原题边界“（2）”：判断并证明积分 \(\int_{1}^{+\infty} \frac{y \sin (x y)}{x\left(1+y^{2}\right)} \mathrm{d} y\) 在 \((0,1)\) 上的一致收敛性．

\subsection{题目4}
原题边界“（3）”：判断并证明函数列 \(f_{n}(x)=n x(1-x)^{n}\) 在 \([0,1]\) 上的一致收敛性．

\subsection{题目6}
六．（20 分）证明：\(\sum_{n=1}^{\infty} \frac{1}{n^{n}}=\int_{0}^{1} x^{-x} \mathrm{~d} x\) ．
\end{document}
